\begin{tcolorbox}[colback=red!3!white,colframe=red!50!black,boxrule=0.8pt,title={Correction notes},fonttitle=\bfseries]
\textbf{Note on the muon $g-2$ status (2025).} This document refers to the deviation between the measured $a_\mu$ and the Standard Model prediction as it stood in 2021. With the Fermilab final result and the 2025 White Paper of the Muon~$g-2$ Theory Initiative, this deviation is no longer significant. Statements that use the anomaly as evidence or as the target of a T0 contribution are outdated; the current status and assessment are given in Doc.~382.
\end{tcolorbox}

\section*{Abstract}
		This work presents a ratio-based formulation of T0 physics that is meant to get by with considerably fewer experimental parameters. Building on the simplified Dirac equation and the universal Lagrangian, we argue that fundamental physics can be described through dimensionless energy scale ratios rather than through assigned parameters. The T0 system requires one SI reference value to connect ratio-based physics to measurable quantities. In natural units, Einstein's $E = mc^2$ becomes the identity $E = m$; the energy relationships that follow hold exactly by definition, while composite multi-parameter formulas agree numerically to 99.98\%. Physics is to be traced back to energy scale ratios, described by the field equation $\partial^2 \Efield = 0$, with quantitative predictions via the single parameter $\xipar$ as SI reference scale (plus integer structural choices as motivated settings).

	\medskip\noindent\textbf{Status markers.} Statements marked \textbf{[S]} are \emph{speculative}: a hint or conjecture that has not yet been derived or numerically checked in the corpus.

	
	
	\section{The T0 Approach: From Parameters to Ratios}
	
	\subsection{The Change of Perspective}
	
	The T0 approach proposes a change of perspective in the understanding of fundamental physics:
	
	\begin{tcolorbox}[colback=red!5!white,colframe=red!75!black,title=Comparison]
		\textbf{Traditional Physics}: Multiple experimental parameters
		\begin{itemize}
			\item $G = 6.67 \times 10^{-11}$ m³/(kg·s²) (measured)
			\item $\alpha = 1/137$ (measured)
			\item $m_e = 9.109 \times 10^{-31}$ kg (measured)
			\item 20+ independent parameters required
		\end{itemize}
		
		\textbf{T0 Ratio-Based Physics}: Dimensionless scale relationships
		\begin{itemize}
			\item Physics through energy scale ratios
			\item One SI reference value for quantitative predictions
			\item Mathematical relationships, not experimental parameters
			\item Pure energy identities: $E = m$, $E = 1/L$, $E = 1/T$
		\end{itemize}
	\end{tcolorbox}
	
	\subsection{Building on T0 Foundations}
	
	This work forms the third stage of a three-stage development:
	
	\textbf{Stage 1 - Simplified Dirac}: Complex 4×4 matrices → Simple field dynamics $\partial^2 \deltam = 0$
	
	\textbf{Stage 2 - Universal Lagrangian}: 20+ fields → One equation $\Lag = \varepsilon \cdot (\partial \deltam)^2$
	
	\textbf{Stage 3 - Ratio-Based Physics}: Multiple parameters → Energy scale ratios + SI reference
	
	\subsection{The Energy Identity}
	
	In natural units ($\hbar = c = 1$), Einstein's equation takes the simple form:
	
	\begin{equation}
		\boxed{E = m}
		\label{eq:energy_mass_identity}
	\end{equation}
	
	In these units it is an \textbf{identity}, not a conversion: mass and energy are treated as the same physical quantity.
	
	\begin{tcolorbox}[colback=blue!5!white,colframe=blue!75!black,title=Universal Energy Relationships]
		\textbf{Energy identity system}:
		\begin{align}
			E &= m \quad \text{(Mass is energy)} \\
			E &= T_{\text{temp}} \quad \text{(Temperature is energy)} \\
			E &= \omega \quad \text{(Frequency is energy)} \\
			E &= \frac{1}{L} \quad \text{(Length is inverse energy)} \\
			E &= \frac{1}{T} \quad \text{(Time is inverse energy)}
		\end{align}
		
		\textbf{Accuracy}: exact by definition (identities in natural units)
		
		\textbf{Complex formulas}: 99.98-100.04\% (rounding errors accumulate)
		
		\textbf{Consequence}: Fewer composite formulas introduce fewer rounding errors.
	\end{tcolorbox}
	
	\section{Part I: Pure Ratio-Based Physics (without SI Reference)}
	
	\subsection{Universal Energy Field Dynamics}
	
	In the T0 framework, all particles are energy excitation patterns in the universal field $\Efield(x,t)$:
	
	\begin{equation}
		\boxed{\partial^2 \Efield = 0}
		\label{eq:universal_field_equation}
	\end{equation}
	
	\textbf{Claim}: This Klein-Gordon equation for energy is meant to describe all particles.
	
	\subsection{Universal Energy Lagrangian}
	
	\begin{equation}
		\boxed{\Lag = \varepsilon \cdot (\partial \Efield)^2}
		\label{eq:universal_lagrangian}
	\end{equation}
	
	where $\varepsilon$ represents the energy scale coupling (dimensionless ratio).
	
	\subsection{Anti-Energy: Symmetry}
	
	\begin{equation}
		\boxed{\Efield_{\text{Antiparticle}} = -\Efield_{\text{Particle}}}
		\label{eq:energy_antisymmetry}
	\end{equation}
	
	\textbf{Physical picture}: Positive and negative energy excitations of the same field.
	
	\textbf{Lagrangian universality}:
	\begin{align}
		\Lag[+\Efield] &= \varepsilon \cdot (\partial \Efield)^2 \\
		\Lag[-\Efield] &= \varepsilon \cdot (\partial \Efield)^2
	\end{align}
	
	Same physics for particles and antiparticles through squaring.
	
	\subsection{Pure Ratio Predictions (Independent of the Value of $\xipar$)}
	
	\subsubsection{Universal Lepton Ratios}
	
	\begin{equation}
		\boxed{\frac{a_e^{(T0)}}{a_{\mu}^{(T0)}} = 1}
		\label{eq:universal_lepton_ratio}
	\end{equation}
	
	\textbf{Physical meaning}: All leptons receive identical energy corrections.
	
	\subsubsection{Energy Independence Ratios}
	
	\begin{equation}
		\boxed{\frac{\Delta\Gamma^{\mu}(E_1)}{\Delta\Gamma^{\mu}(E_2)} = 1}
		\label{eq:energy_independence_ratio}
	\end{equation}
	
	\textbf{Distinguishing feature}: In contrast to Standard Model running couplings.
	
	\section{Part II: Quantitative Predictions (SI Reference Required)}
	
	\subsection{The SI Reference Scale}
	
	To make quantitative predictions, T0 physics needs a connection to the SI system:
	
	\begin{tcolorbox}[colback=green!5!white,colframe=green!75!black,title=SI Reference Scale]
		\textbf{Definition}: $\xipar$ is a dimensionless energy scale ratio and the only parameter of the framework; it is geometrically motivated, not fitted to measured data.
		
		\textbf{Higgs energy ratio}:
		\begin{equation}
			\xipar = \frac{\lambda_h^2 v^2}{16\pi^3 E_h^2}
		\end{equation}
		
		\textbf{Geometric energy ratio}:
		\begin{equation}
			\xipar = \frac{2\ell_P}{\lambda_C}
		\end{equation}
		
		\textbf{SI reference value}: $\xipar = 1.33 \times 10^{-4}$
		
		\textbf{Role}: Connects dimensionless ratios to SI-measurable quantities
	\end{tcolorbox}
	
	\subsection{Quantitative Lepton Predictions}
	
	With the SI reference scale:
	
	\begin{equation}
		a_{\ell}^{(T0)} = \frac{1}{2\pi} \times \xipar^2 \times \frac{1}{12}
		\label{eq:quantitative_lepton_correction}
	\end{equation}
	
	\textbf{Numerical calculation}:
	\begin{align}
		a_{\ell}^{(T0)} &= \frac{1}{2\pi} \times (1.33 \times 10^{-4})^2 \times \frac{1}{12} \\
		&= \frac{1}{6.283} \times 1.77 \times 10^{-8} \times 0.0833 \\
		&= 2.47 \times 10^{-10}
	\end{align}
	
	\begin{tcolorbox}[colback=blue!5!white,colframe=blue!75!black,title=Universal Lepton Prediction]
		\textbf{Electron g-2}: $a_e^{(T0)} = 2.47 \times 10^{-10}$
		
		\textbf{Muon g-2}: $a_{\mu}^{(T0)} = 2.47 \times 10^{-10}$ (identical)
		
		\textbf{Tau g-2}: $a_{\tau}^{(T0)} = 2.47 \times 10^{-10}$ (same value)
		
		\textbf{Current muon anomaly}: $\Delta a_{\mu} \approx 25 \times 10^{-10}$
		
		\textbf{T0 contribution}: $\sim 10\%$ of the observed anomaly
	\end{tcolorbox}
	
	\subsection{Quantitative QED Predictions}
	
	\begin{equation}
		\frac{\Delta\Gamma^{\mu}}{\Gamma^{\mu}} = \xipar^2 = 1.77 \times 10^{-8}
		\label{eq:quantitative_qed_correction}
	\end{equation}
	
	\textbf{Energy independence verification}:
	\begin{table}[htbp]
		\centering
		\adjustbox{max width=\textwidth}{%
\begin{tabular}{lcc}
			\toprule
			\textbf{Energy Scale} & \textbf{T0 Correction} & \textbf{Standard Model} \\
			\midrule
			1 MeV & $1.77 \times 10^{-8}$ & Running $\alpha(E)$ \\
			1 GeV & $1.77 \times 10^{-8}$ & Running $\alpha(E)$ \\
			100 GeV & $1.77 \times 10^{-8}$ & Running $\alpha(E)$ \\
			1 TeV & $1.77 \times 10^{-8}$ & Running $\alpha(E)$ \\
			\bottomrule
		\end{tabular}%
}
		\caption{Energy-independent T0 corrections vs. Standard Model}
	\end{table}
	
	\section{Experimental Verification Strategy}
	
	\subsection{Pure Ratio Tests (No SI Reference Needed)}
	
	\textbf{Test 1 - Universal lepton ratios}:
	\begin{itemize}
		\item Measure $a_e^{(T0)}/a_{\mu}^{(T0)} = 1$
		\item Independent of absolute values
		\item Directly tests universality principle
	\end{itemize}
	
	\textbf{Test 2 - Energy independence}:
	\begin{itemize}
		\item Measure QED corrections at different energies
		\item Ratio should be constant: $\Delta\Gamma(E_1)/\Delta\Gamma(E_2) = 1$
		\item Distinguishes from Standard Model running couplings
	\end{itemize}
	
	\textbf{Test 3 - Wavelength ratios} \textbf{[S]}:
	\begin{itemize}
		\item Multi-wavelength observations of same objects
		\item Test $z(\lambda_1)/z(\lambda_2) = \lambda_2/\lambda_1$
		\item Independent of absolute redshift calibration
	\end{itemize}
	
	\subsection{Quantitative Tests (Require SI Reference)}
	
	\textbf{Precision g-2 measurements}:
	\begin{itemize}
		\item Electron g-2: Detect $2.47 \times 10^{-10}$ correction
		\item Muon g-2: Test $\sim 10\%$ of current anomaly
		\item Tau g-2: First measurement, expect same value
	\end{itemize}
	
	\textbf{Multi-energy QED tests}:
	\begin{itemize}
		\item Measure absolute $\Delta\Gamma/\Gamma = 1.77 \times 10^{-8}$
		\item Verify energy independence across decades
		\item Compare with Standard Model predictions
	\end{itemize}
	
	\section{Dark Matter and Dark Energy\\ from Energy Ratios}
	
	\textbf{[S]} The statements in this section are speculative: the present corpus deliberately leaves the cosmic sector aside, because the Standard Model does not derive it either (P39).
	
	\subsection{Dark Matter: Sub-threshold Energy Oscillations}
	
	\textbf{Ratio-based description}:
	\begin{equation}
		\frac{\Efield_{\text{dark}}}{\Efield_{\text{threshold}}} = \xipar \sqrt{\frac{\rho_{\text{local}}}{\rho_{\text{critical}}}}
	\end{equation}
	
	\textbf{Physical mechanism}: Random phase energy oscillations below particle detection threshold.
	
	\subsection{Dark Energy: Large-scale Energy Gradients}
	
	\textbf{Ratio-based energy density}:
	\begin{equation}
		\frac{\rho_{\Lambda}}{\rho_{\text{critical}}} = \frac{1}{2} \xipar^2 \left(\frac{E_{\text{Planck}}}{L_{\text{Hubble}} \cdot E_{\text{Planck}}}\right)^2
	\end{equation}
	
	\textbf{Quantitative prediction}: $\rho_{\Lambda} \approx 6 \times 10^{-30}$ g/cm$^3$ (of the order of the observed value)
	
	\section{Philosophical Classification: Physics without Independent Matter}
	
	\subsection{Pure Energy Reality}
	
	\begin{tcolorbox}[colback=purple!5!white,colframe=purple!75!black,title=Dematerialized Description]
		\textbf{Traditional view}: Matter, energy, forces, spacetime as separate entities
		
		\textbf{T0 view}: Only energy patterns and their ratios
		
		\textbf{What we call particles}: Localized energy concentrations
		
		\textbf{What we call forces}: Energy gradient interactions
		
		\textbf{What we call spacetime}: Energy pattern substrate
		
		\textbf{What we call consciousness} \textbf{[S]}: Self-referential energy patterns
		
		\textbf{Core statement}: Pure energy relationships governed by $\partial^2 \Efield = 0$
	\end{tcolorbox}
	
	\subsection{From Complexity to Simplicity}
	
	\textbf{Physics evolution}:
	\begin{enumerate}
		\item \textbf{Antiquity}: Four elements
		\item \textbf{Classical}: Particles in spacetime
		\item \textbf{Modern}: Fields and forces
		\item \textbf{Standard Model}: 20+ parameters
		\item \textbf{T0 approach}: Energy ratios + one SI reference
	\end{enumerate}
	
	\textbf{The aim is far-reaching simplification}: as few fundamental assumptions as possible.
	
	\subsection{Consciousness and Energy Patterns}
	
	\textbf{A further question}: If everything is energy patterns, what about consciousness?
	
	\textbf{[S] T0 interpretation}: Consciousness would be a self-observing energy pattern. We are temporary organizations of the universal energy field that have developed the ability for self-reference and subjective experience.
	
	\section{The Ratio Physics Legacy}
	
	\subsection{Intended Results}
	
	The ratio-based T0 approach pursues the following goals; the respective status is recorded in the detailed documents:
	
	\begin{enumerate}
		\item \textbf{Parameters reduced}: 20+ → $\xipar$ as SI reference (plus integer settings)
		\item \textbf{Forces described together}: Through energy gradient interactions
		\item \textbf{Particle variety ordered}: Described as energy patterns
		\item \textbf{Antiparticles interpreted}: Negative energy excitations
		\item \textbf{Gravitation included}: Via energy-spacetime coupling
		\item \textbf{Dark phenomena} \textbf{[S]}: Interpreted as energy field effects (speculative, P39)
		\item \textbf{Exact identities}: By definition in natural units
		\item \textbf{Ratio-based physics formulated}: Pure scale relationships
	\end{enumerate}
	
	\subsection{The Two-Stage Testing Strategy}
	
	\textbf{Stage 1 - Pure ratios} (independent of the value of $\xipar$):
	\begin{itemize}
		\item Universal lepton correction ratios
		\item Energy-independent QED ratios
		\item Wavelength-dependent redshift ratios \textbf{[S]}
		\item Gravitational modification ratios
	\end{itemize}
	
	\textbf{Stage 2 - Quantitative predictions} (SI reference):
	\begin{itemize}
		\item Absolute g-2 corrections
		\item Absolute QED vertex modifications
		\item Absolute cosmological parameters \textbf{[S]}
		\item Absolute dark matter/energy densities \textbf{[S]}
	\end{itemize}
	
	\subsection{Status of the Framework}
	
	\begin{tcolorbox}[colback=yellow!5!white,colframe=orange!75!black,title=Basic Structure in the T0 Framework]
		\textbf{In the T0 framework, the basic structure reduces to a few building blocks}.
		
		\textbf{The fundamental equation}: $\partial^2 \Efield = 0$
		
		\textbf{The universal ratios}: Energy scale relationships
		
		\textbf{The SI connection}: One reference scale $\xipar$
		
		\textbf{Everything else}: Various solutions and patterns
		
		\textbf{Deeper level}: Whether a deeper level exists remains open; the T0 framework does not provide for one
		
		\textbf{Future work}: Mainly applications and measurements that test the framework
	\end{tcolorbox}
	
	\section{Conclusion: The Ratio-Based Universe}
	
	\subsection{Summary of the Structure}
	
	In the T0 framework, physics is described through energy scale ratios:
	
	\textbf{Level 1}: Dimensionless energy ratios (independent of the SI reference)
	
	\textbf{Level 2}: One SI reference scale (quantitative predictions)
	
	\textbf{Level 3}: Pure energy patterns governed by $\partial^2 \Efield = 0$
	
	The observed and measured phenomena are meant to follow from this simple ratio-based structure.
	
	\subsection{Review}
	
	The path leads from the many parameters of traditional physics to a simpler, ratio-based energy dynamics.
	
	\textbf{The core statement}: In this view, the core of nature lies not in complicated mathematics or exotic phenomena, but in simple scale relationships.
	
	\textbf{One field}. \textbf{One equation}. \textbf{Energy ratios}. \textbf{One SI reference}.
	
	Everything further would then be various patterns and ratios of energy, up to the pattern we call human consciousness \textbf{[S]}.
	
	\begin{equation}
		\boxed{\text{Reality} = \text{Energy ratios in } \Efield(x,t)}
	\end{equation}
	
	\textbf{In summary: In the T0 framework, the universe is described as a structure of energy ratios; this is tested through the tests listed above.}
	
\begin{thebibliography}{99}

\bibitem{pascher_simplified_dirac_2025}
Pascher, J. (2025). \textit{Vereinfachte Dirac-Gleichung in der T0-Theorie: Von komplexen 4×4-Matrizen zu einfacher Feld-Knoten-Dynamik}. \\
\href{https://github.com/jpascher/T0-Time-Mass-Duality/blob/main/2/pdf/050_diracVereinfacht_En.pdf}{https://github.com/jpascher/T0-Time-Mass-Duality/blob/main/2/pdf/\\050\_diracVereinfacht\_En.pdf}

\bibitem{pascher_lagrangian_comparison_2025}
Pascher, J. (2025). \textit{Einfache Lagrange-Revolution: Von Standardmodell-Komplexität zu T0-Eleganz}. \\
\href{https://github.com/jpascher/T0-Time-Mass-Duality/blob/main/2/pdf/049_LagrandianVergleich_En.pdf}{https://github.com/jpascher/T0-Time-Mass-Duality/blob/main/2/pdf/\\049\_LagrandianVergleich\_En.pdf}

\bibitem{pascher_verification_table_2025}
Pascher, J. (2025). \textit{T0-Modell-Verifikation: Skalen-Verhältnis-basierte Berechnungen vs. CODATA/Experimentelle Werte}. \\
\href{https://github.com/jpascher/T0-Time-Mass-Duality/blob/main/2/pdf/054_Elimination_Of_Mass_Dirac_Tabelle_En.pdf}{https://github.com/jpascher/T0-Time-Mass-Duality/blob/main/2/pdf/\\054\_Elimination\_Of\_Mass\_Dirac\_Tabelle\_En.pdf}

\bibitem{einstein_mass_energy_1905}
Einstein, A. (1905). \textit{Ist die Trägheit eines Körpers von seinem Energieinhalt abhängig?} Ann. Phys. \textbf{18}, 639--641.

\bibitem{dirac_original_1928}
Dirac, P. A. M. (1928). \textit{The Quantum Theory of the Electron}. Proc. R. Soc. London A \textbf{117}, 610.

\bibitem{muong2_experiment_2021}
Myon g-2 Kollaboration (2021). \textit{Messung des positiven Myon-anomalen magnetischen Moments auf 0.46 ppm}. Phys. Rev. Lett. \textbf{126}, 141801.

\bibitem{higgs_mechanism_1964}
Higgs, P. W. (1964). \textit{Gebrochene Symmetrien und die Massen von Eichbosonen}. Phys. Rev. Lett. \textbf{13}, 508--509.

\bibitem{planck_collaboration_2020}
Planck Kollaboration (2020). \textit{Planck 2018 Ergebnisse. VI. Kosmologische Parameter}. Astron. Astrophys. \textbf{641}, A6.

\bibitem{weinberg_qft_1995}
Weinberg, S. (1995). \textit{Die Quantentheorie der Felder, Band 1: Grundlagen}. Cambridge University Press.

\bibitem{particle_data_group_2022}
Teilchendaten-Gruppe (2022). \textit{Übersicht der Teilchenphysik}. Prog. Theor. Exp. Phys. \textbf{2022}, 083C01.

\end{thebibliography}
	